\begin{lemma}
\label{lemma-derived-of-quotient}
Let $\mathcal{A}$ be an abelian category.
Let $\mathcal{B} \subset \mathcal{A}$ be a Serre subcategory.
Then $D(\mathcal{A}) \to D(\mathcal{A}/\mathcal{B})$
is essentially surjective.
\end{lemma}

\begin{proof}
We will use the description of the category $\mathcal{A}/\mathcal{B}$
in the proof of
Homology, Lemma \ref{homology-lemma-serre-subcategory-is-kernel}.
Let $(X^\bullet, d^\bullet)$ be a complex of $\mathcal{A}/\mathcal{B}$.
This means that $X^i$ is an object of $\mathcal{A}$ and
$d^i : X^i \to X^{i + 1}$ is a morphism in $\mathcal{A}/\mathcal{B}$
such that $d^i \circ d^{i - 1} = 0$ in $\mathcal{A}/\mathcal{B}$.

\medskip\noindent
For $i \geq 0$ we may write $d^i = (s^i, f^i)$ where $s^i : Y^i \to X^i$
is a morphism of $\mathcal{A}$ whose kernel and cokernel are in $\mathcal{B}$
(equivalently $s^i$ becomes an isomorphism in the quotient category)
and $f^i : Y^i \to X^{i + 1}$ is a morphism of $\mathcal{A}$.
By induction we will construct a commutative diagram
$$
\xymatrix{
& (X')^1 \ar@{..>}[r] & (X')^2 \ar@{..>}[r] & \ldots \\
X^0 \ar@{..>}[ru] &
X^1 \ar@{..>}[u] &
X^2 \ar@{..>}[u] &
\ldots \\
Y^0 \ar[u]_{s^0} \ar[ru]_{f^0} &
Y^1 \ar[u]_{s^1} \ar[ru]_{f^1} &
Y^2 \ar[u]_{s^2} \ar[ru]_{f^2} &
\ldots
}
$$
where the vertical arrows $X^i \to (X')^i$ become isomorphisms
in the quotient category. Namely, we first let
$(X')^1 = \Coker(Y^0 \to X^0 \oplus X^1)$ (or rather the
pushout of the diagram with arrows $s^0$ and $f^0$) which gives the
first commutative diagram. Next, we take
$(X')^2 = \Coker(Y^1 \to (X')^1 \oplus X^2)$. And so on.
Setting additionally $(X')^n = X^n$ for $n \leq 0$ we see that the map
$(X^\bullet, d^\bullet) \to ((X')^\bullet, (d')^\bullet)$
is an isomorphism of complexes in $\mathcal{A}/\mathcal{B}$.
Hence we may assume $d^n : X^n \to X^{n + 1}$ is given
by a map $X^n \to X^{n + 1}$ in $\mathcal{A}$ for $n \geq 0$.

\medskip\noindent
Dually, for $i < 0$ we may write $d^i = (g^i, t^{i + 1})$ where
$t^{i + 1} : X^{i + 1} \to Z^{i + 1}$ is an isomorphism in the
quotient category and $g^i : X^i \to Z^{i + 1}$ is a morphism.
By induction we will construct a commutative diagram
$$
\xymatrix{
\ldots &
Z^{-2} &
Z^{-1} &
Z^0 \\
\ldots &
X^{-2} \ar[u]_{t_{-2}} \ar[ru]_{g_{-2}} &
X^{-1} \ar[u]_{t_{-1}} \ar[ru]_{g_{-1}} &
X^0 \ar[u]_{t^0} \\
\ldots &
(X')^{-2} \ar@{..>}[u] \ar@{..>}[r] &
(X')^{-1} \ar@{..>}[u] \ar@{..>}[ru]
}
$$
where the vertical arrows $(X')^i \to X^i$ become isomorphisms
in the quotient category. Namely, we take
$(X')^{-1} = X^{-1} \times_{Z^0} X^0$. Then we take
$(X')^{-2} = X^{-2} \times_{Z^{-1}} (X')^{-1}$. And so on.
Setting additionally $(X')^n = X^n$ for $n \geq 0$ we see that the map
$((X')^\bullet, (d')^\bullet) \to (X^\bullet, d^\bullet)$
is an isomorphism of complexes in $\mathcal{A}/\mathcal{B}$.
Hence we may assume $d^n : X^n \to X^{n + 1}$ is given
by a map $d^n : X^n \to X^{n + 1}$ in $\mathcal{A}$
for all $n \in \mathbf{Z}$.

\medskip\noindent
In this case we know the compositions $d^n \circ d^{n - 1}$
are zero in $\mathcal{A}/\mathcal{B}$. If for $n > 0$ we replace
$X^n$ by
$$
(X')^n = X^n/\sum\nolimits_{0 < k \leq n} \Im(\Im(X^{k - 2} \to X^k) \to X^n)
$$
then the compositions $d^n \circ d^{n - 1}$ are zero for $n \geq 0$.
(Similarly to the second paragraph above we obtain an isomorphism of
complexes
$(X^\bullet, d^\bullet) \to ((X')^\bullet, (d')^\bullet)$.)
Finally, for $n < 0$ we replace $X^n$ by
$$
(X')^n = \bigcap\nolimits_{n \leq k < 0}
(X^n \to X^k)^{-1}\Ker(X^k \to X^{k + 2})
$$
and we argue in the same manner to get a complex in $\mathcal{A}$
whose image in $\mathcal{A}/\mathcal{B}$ is isomorphic to the given one.
\end{proof}